%!TEX root = ../thesis.tex
% Chapter --- Conclusion

\chapter{Conclusion} \label{sec:conclusion}

This report has presented a Python replication of the wavelet scattering, CNN, and CNN~+~WS fusion branches of Razali et al.\ \cite{razali2023cnn}, evaluated on both the fatty versus fibroglandular tissue task and the benign versus malignant mass task, together with two extensions to the wavelet-scattering branch: a supervisor-suggested scattering covariance representation~\cite{cheng2021scattering} for the tissue task, and a region-aware wavelet-scattering decomposition developed for the mass task. The work proceeded iteratively: an early CNN run on the tissue task surfaced a data-leakage issue that prompted a re-examination of the train--test split methodology for both branches; only after that correction was the mass-classification task added, followed by the fusion model and finally the scattering covariance extension. The chronological development log is in Section~\ref{sec:results}; Tables~\ref{tab:results_tissue} and~\ref{tab:results_mass} compare every replicated configuration against the paper's reported benchmarks.

Under the final five-seed image-level protocol (Tables~\ref{tab:results_tissue}--\ref{tab:results_mass}), all three paper configurations reproduce for both tasks. On \emph{tissue}, WS-only reaches AU-ROC $0.895 \pm 0.070$ (test $0.839 \pm 0.071$), trailing the paper by roughly $7$--$8$~pp --- the largest replication gap, attributable to standard CLAHE versus the paper's density-aware contrast and to scattering-library conventions --- while the CNN-only ($0.979 \pm 0.010$ AU-ROC, test $0.979 \pm 0.010$) and CNN~+~WS fusion ($0.988 \pm 0.007$, test $0.964 \pm 0.030$) are near-saturated and slightly exceed the paper, reflecting a cleaner preprocessing pipeline rather than a methodological change. On \emph{mass}, WS-only is weak (AU-ROC $0.726 \pm 0.030$, test $0.701 \pm 0.075$), the CNN-only model reproduces the paper's reported values to within a few points on the seed-34 metrics (including the overfitting signature in its training curves), and CNN~+~WS fusion is the strongest strict replication model (AU-ROC $0.816 \pm 0.057$, test $0.792 \pm 0.058$) --- above both CNN-only ($0.762 \pm 0.117$) and WS-only ($0.726$) --- directly reproducing the paper's central narrative that handcrafted scattering features compensate for small-sample CNN overfitting.

The scattering covariance extension (Section~\ref{sec:scov_tissue}) was tested against the original WS baseline, a direct usage of the Cheng package, and the paper's published WS-only number. The final adopted WS-branch representation is the adapted scattering covariance configuration with $J=6, L=3$ on the \texttt{for\_synthesis\_iso} vector ($580$ scalar coefficients per patch), selected by a small $J$/$L$ grid search ranked by 10-fold group-aware CV mean on the training partition. This configuration reaches AU-ROC $0.984 \pm 0.010$ / test $0.959 \pm 0.013$ on tissue (versus averaged WS's $0.895$/$0.839$), the strongest adopted WS-branch tissue result in the project. On \emph{mass}, however, it gives no meaningful AU-ROC gain (J6/L3 scattering covariance $0.734 \pm 0.062$ versus averaged WS $0.726 \pm 0.030$), so it is adopted as a tissue-side extension only; the mass improvement comes instead from preserving the scattering spatial map and learning over it (Section~\ref{sec:region_aware}). The scattering covariance representation respects the supervisor's no-flattening constraint by construction: its outputs are scalar coefficients indexed by scale and orientation tuples, never a stack of flattened spatial maps. The original Mallat-style WS implementation is retained as the faithful paper-replication reference; the Direct Cheng usage at $L=4$ is retained as an independent sensitivity check confirming the result is not specific to the $L=3$ choice. The CNN+WS fusion experiments (Section~\ref{sec:stage6}) used the original $S_0, S_1, S_2$ WS features for direct comparability with the paper; reproducing fusion on top of the scattering covariance features is identified as natural further work.

The GLCM branch (Models~3, 4, 6 in the paper's Table~6) was excluded on supervisor advice; the paper itself shows GLCM contributes negligible improvement when combined with WS, so the practical loss is small. A small data-quality issue was uncovered during the scattering covariance experiment: one mass annotation (file\_id 50994354, Section~\ref{sec:mass_quality_filter}) had a bounding box outside the breast mask, producing a uniformly-zero patch. A defensive filter was added to every mass-using notebook; the test set went from 25 to 24 mass patches with no qualitative change to method rankings.

Beyond the tissue-side scattering covariance, the project's main \emph{mass}-side extension is to preserve the wavelet-scattering spatial map and learn over it with a small CNN (Section~\ref{sec:region_aware}). Because averaged WS discards \emph{where} lesion structure sits, the maps are retained at $J=3$ rather than collapsed to a scalar and a small WS-map CNN learns over them; in the lesion-relative \emph{region-aware} instantiation each lesion is additionally decomposed into core/inner, boundary, whole-lesion, and surrounding-context regions before scattering. This recovers most of the lost ground: a whole-patch $J=3$-map CNN reaches AU-ROC $0.909 \pm 0.053$, the region-aware classical (logistic-regression) model reaches $0.892 \pm 0.041$, and the region-aware WS$\rightarrow$CNN head reaches $0.946 \pm 0.036$ on the five-seed holdout (corroborated by a tight group-aware cross-validation of $0.938 \pm 0.007$) --- all well above both the $0.726$ averaged-WS baseline and the $0.816$ paper-style fusion. A region ablation shows the all-four-region tensor is the strongest complete configuration, while neighbouring strong variants overlap within seed noise, so the region split is retained as a principled, interpretable lesion representation and as the current strongest complete region-aware result. A decision-level \emph{late fusion} of the region-aware head with a frozen-ResNet improved-CNN comparator ($0.857 \pm 0.059$) is the numerically strongest configuration ($0.948 \pm 0.033$ AU-ROC, test $0.883$), but it is an ensemble of two models rather than a new representation, and its margin over the region-aware head alone is within seed noise; it is therefore reported as the best \emph{number}, not the central contribution.

The methodological takeaway is twofold. For the \emph{tissue} (texture) task, a richer \emph{scalar} wavelet-scattering descriptor --- scattering covariance --- is effective and sufficient. For the \emph{mass} task, by contrast, \emph{preserving the spatial structure} of the scattering representation --- retaining its lower-scale maps and learning over them with a CNN --- matters more than any change of scalar feature or classifier; conditioning those maps on lesion morphology (core, margin, surrounding context) makes that spatial representation interpretable and, in the current seed-99 results, numerically strongest among the complete region-aware configurations. This is the project's principal finding beyond replication.

These conclusions are bounded by the usual caveats. The INbreast mass set is small ($115$ patches, ${\sim}24$ held out per split), so every mass estimate carries an across-seed band and is reported as a five-seed mean~$\pm$~standard deviation under strict image-level (group-aware) splitting to avoid the patch-level leakage diagnosed early in the project; no data augmentation is used, following the paper; and only a single dataset is evaluated, leaving cross-dataset generalisation to future work. The region-aware WS$\rightarrow$CNN result is reported as a five-seed $0.946$ holdout, corroborated by the $0.938$ group-aware cross-validation.


% ============================================================
% REFERENCES
% ============================================================
